% figbank.tex -- auto-generated figure bank (FLITS scattering/scintillation campaign).
% Drop-in figures with labels + draft captions; move each \begin{figure} into the
% relevant section (results/discussion) and replace the citation placeholders with
% \citep{} keys from bib/refs. Regenerated by analysis nladder/figbank.py.

\begin{figure}
  \centering
  \includegraphics[width=0.95\columnwidth]{pbf_shapes}
  \caption{Pulse-broadening function (PBF) shapes used in the joint fit. A Gaussian
  intrinsic pulse convolved with a thin-screen exponential PBF (green) versus
  power-law-spectrum PBFs with density index $\beta=3.00$ and $\beta=11/3$
  (Kolmogorov). Smaller $\beta$ yields a heavier $t^{-\beta/2}$ tail; $\beta\to4$
  recovers the exponential. [cite: power-law-spectrum PBF, Cordes review \S11.2]}
  \label{fig:pbf_shapes}
\end{figure}

\begin{figure*}
  \centering
  \includegraphics[width=0.95\textwidth]{wilhelm_pbf_evidence}
  \caption{Band-dependent PBF preference for FRB~wilhelm. \emph{Left:} change in
  log-evidence relative to the exponential PBF versus power-law index $\beta$ (best
  single-component, per band). CHIME (400--800\,MHz) mildly prefers a Kolmogorov
  tail ($\Delta\ln Z=+6.2$ at $\beta=11/3$); DSA (1.2--1.5\,GHz) rejects
  heavier-than-exponential tails monotonically. \emph{Right:} joint CHIME+DSA
  evidence (C1D1). A \emph{per-band} PBF (CHIME power-law, DSA exponential) beats the
  all-exponential model by $\Delta\ln Z=+4.0$ ($\sim6\sigma$), whereas a global
  power-law loses $84$. The two bands require different PBFs.}
  \label{fig:pbf_evidence}
\end{figure*}

\begin{figure}
  \centering
  \includegraphics[width=0.95\columnwidth]{wilhelm_scint_dnud_ne2025}
  \caption{Scintillation decorrelation bandwidth $\Delta\nu_d$ versus frequency for
  FRB~wilhelm. DSA-resolved $\Delta\nu_d\approx0.13$\,MHz (blue; nearly flat,
  $x=-0.23\pm0.19$ over the 10\% DSA lever arm); CHIME upper limits (red) since the
  diffractive scale is sub-channel at native resolution. Purple stars: the NE2025
  Milky-Way scattering floor, $\sim8\times$ \emph{above} the measurement. Green
  dash-dot: the pulse-broadening screen ($2\pi\tau\,\Delta\nu_d=C_1=1$), $\sim85\times$
  \emph{below}. The three separated scales imply the dominant pulse-broadening screen
  is neither the Galactic floor nor the resolved scintillation screen (a strong
  host/intervening screen). [cite: NE2025 Galactic electron-density model]}
  \label{fig:scint_screens}
\end{figure}

\begin{figure}
  \centering
  \includegraphics[width=0.95\columnwidth]{codetection_scint_excess}
  \caption{Is the wilhelm scattering excess sightline-specific or systematic?
  Scattering excess (NE2025 Milky-Way-floor $\Delta\nu_d$ divided by the measured
  DSA diffractive $\Delta\nu_d$; excess${=}1$ is the MW floor, dashed) versus
  Galactic latitude $|b|$ for the 12 CHIME--DSA co-detections. Error bars combine
  the measurement error with NE2025's floor uncertainty ($\sigma\approx0.4$\,dex,
  $\sim\!\times2.5$); the annotated $z$ is the significance above the MW floor.
  $\Delta\nu_d$ is the narrowest consistent cluster of well-constrained
  Lorentzian/Gen-Lorentz widths across all stored per-subband scintillation fits
  (rejecting fits railed at the $0.060$\,MHz bound). Filled circles: cleanly-resolved
  diffractive scales significant at $z>2$ -- zach ($10.7\times$), wilhelm
  ($8.0\times$), hamilton ($7.2\times$), all in the shaded band. Open circle: casey
  ($9.6\times$, single-subband, low confidence). Light-blue: clean but $z<2$
  (chromatica $5.9\times$, oran $2.6\times$). Orange up-triangles: lower limits
  (single-component or rail-only fits). The grey band ($|b|<14^\circ$) marks
  sightlines where NE2025 intersects voids and the floor is unreliable (johndoeii,
  isha, mahi fall below unity there). The excess clusters at $\sim7$--$11\times$ on
  $3$--$4$ independent mid-$|b|$ sightlines where the floor is reliable and the
  diffractive scale resolvable (three high-confidence: zach, wilhelm, hamilton;
  casey adds a fourth at high $|b|$ but is low-confidence), and never the reverse --
  evidence for a common (host/intervening) scattering enhancement rather than a
  wilhelm-only peculiarity. Per-sightline significance is $\sim\!2\sigma$, dominated
  by NE2025's floor uncertainty; because that error is a shared systematic across
  sightlines (not independent), the result rests on the \emph{asymmetry} -- the
  measurement is always in excess where the floor is reliable, never below -- rather
  than on a stacked significance. [cite: NE2025 Galactic electron-density model]}
  \label{fig:codetection_excess}
\end{figure}
